% ============================================================
% Sección 4: Resultados
% ============================================================

\section{Resultados}
\label{sec:resultados}

\subsection{Validación 1D --- Helmholtz con solución analítica}
\label{subsec:resultados1d}

La \cref{tab:resultados1d} muestra las métricas de precisión obtenidas para la
ecuación de Helmholtz 1D con la arquitectura SIREN $5\times64$, $\omega_0 = 1.0$.

\begin{table}[htbp]
\centering
\begin{threeparttable}
\caption{Métricas de precisión --- PINN-SIREN 1D ($k = 2\pi$)}
\label{tab:resultados1d}
\begin{tabular}{lc}
\toprule
Métrica & Valor \\
\midrule
Error L2 (\%)          & 0.006 \\
MSE                    & $1.5 \times 10^{-9}$ \\
RMSE                   & $3.9 \times 10^{-5}$ \\
MAE                    & $3.50 \times 10^{-5}$ \\
Error máximo           & $5.77 \times 10^{-5}$ \\
$R^2$                  & 1.000000 \\
Correlación de Pearson & 1.000000 \\
\midrule
Épocas Adam            & 15{,}000 / 15{,}000 \\
Iteraciones L-BFGS     & 202 / 500 \\
Tiempo total           & 251 s \\
\bottomrule
\end{tabular}
\begin{tablenotes}
\small
\item Evaluado en malla uniforme de 1{,}000 puntos en $[0,1]$.
      Arquitectura SIREN: 5 capas $\times$ 64 neuronas, $\omega_0 = 1.0$.
      Solución analítica: $E_{\mathrm{exacta}}(x) = \cos(kx)$.
      GPU: NVIDIA RTX 5050, SEED = 42.
\end{tablenotes}
\end{threeparttable}
\end{table}

El error L2 de $0.006\%$ supera el umbral de la tesis ($< 5\%$) por casi tres órdenes de magnitud
y mejora en un factor de $\times 415$ el resultado de \citet{Schoder2024_PINN_acoustics},
que reporta un error del $2.49\%$ para Helmholtz 1D con PINNs convencionales.
El ajuste es prácticamente perfecto: $R^2 = 1.000000$ y correlación de Pearson $= 1.000000$
en toda la malla de evaluación.

\begin{figure}[htbp]
  \centering
  \includegraphics[width=0.9\linewidth]{resultados_pinn_1d.png}
  \caption{Validación PINN-SIREN para Helmholtz 1D.
           (a)~Comparación entre la predicción de la red y la solución analítica
           $E_{\mathrm{exacta}}(x) = \cos(kx)$ en 1{,}000 puntos uniformes en $[0,1]$.
           (b)~Error puntual $|\hat{E}(x) - E_{\mathrm{exacta}}(x)|$.
           Error $L^2 = 0.006\%$, $R^2 = 1.000000$.
           GPU: NVIDIA RTX 5050, SEED = 42.}
  \label{fig:resultados1d}
\end{figure}

\begin{figure}[htbp]
  \centering
  \includegraphics[width=\linewidth]{metricas_adicionales_1d.png}
  \caption{Métricas adicionales de entrenamiento NB01.
           De izquierda a derecha: curva de pérdida total (Adam + L-BFGS),
           distribución del error puntual y métricas de ajuste ($R^2$, Pearson).
           La transición Adam $\to$ L-BFGS a la época 15{,}000 es visible como
           discontinuidad en el eje de iteraciones.}
  \label{fig:metricas1d}
\end{figure}

\subsection{Validación 2D --- Helmholtz con campo complejo}
\label{subsec:resultados2d}

La \cref{tab:resultados2d} reporta las métricas para la validación 2D con la arquitectura
optimizada SIREN $5\times128$, $\omega_0 = 1.0$, y muestreo LHS de 3{,}000 puntos.

\begin{table}[htbp]
\centering
\begin{threeparttable}
\caption{Métricas de precisión --- PINN-SIREN 2D ($k = 2\pi$, $k_x = k_y = k/\sqrt{2}$)}
\label{tab:resultados2d}
\begin{tabular}{lcc}
\toprule
Métrica & $E_{\mathrm{real}}$ & $E_{\mathrm{imag}}$ \\
\midrule
Error L2 (\%)          & 0.214 & 0.127 \\
MSE                    & $2.20 \times 10^{-6}$ & $8.43 \times 10^{-7}$ \\
RMSE                   & $1.48 \times 10^{-3}$ & $9.18 \times 10^{-4}$ \\
MAE                    & $1.27 \times 10^{-3}$ & $7.71 \times 10^{-4}$ \\
Error máximo           & $3.05 \times 10^{-3}$ & $2.93 \times 10^{-3}$ \\
$R^2$                  & 0.999995 & 0.999998 \\
Correlación de Pearson & 0.999998 & 0.999999 \\
\midrule
Error L2 promedio (\%) & \multicolumn{2}{c}{0.171} \\
\bottomrule
\end{tabular}
\begin{tablenotes}
\small
\item Evaluado en malla uniforme de $100 \times 100$ puntos.
      Arquitectura: SIREN $5 \times 128$, $\omega_0 = 1.0$, 66{,}690 parámetros.
      Puntos de colocación: LHS $N_c = 3{,}000$.
      Frontera: onda plana $e^{i(k_x x + k_y y)}$ en 4 bordes, $N_b = 300$ pts/borde.
      Entrenamiento: Adam (8{,}737 épocas) + L-BFGS (1{,}035 iter), tiempo total 299 s.
      Varianza entre semillas (SEED $\in \{42, 123, 777\}$):
      error $L^2$ promedio $= 0.155 \pm 0.020$\,\%.
\end{tablenotes}
\end{threeparttable}
\end{table}

El error L2 promedio de $0.171\%$ supera el umbral de la tesis ($< 5\%$) por más de un
orden de magnitud y mejora en $\times 11.2$ el resultado de
\citet{Schoder2024_PINN_acoustics} para Helmholtz 2D ($1.91\%$).
La precisión se mantiene alta en ambas componentes del campo complejo, con $R^2 > 0.9999$
para $E_{\mathrm{real}}$ y $E_{\mathrm{imag}}$.

El aumento de capacidad de $d = 64$ a $d = 128$ fue determinante: con $d = 64$ el error
promedio era $0.436\%$; con $d = 128$ se redujo a $0.171\%$, una mejora de $2.5\times$.
El muestreo LHS contribuyó a la reducción del error en regiones de baja densidad de
puntos respecto a mallas cartesianas equivalentes.

\begin{figure}[htbp]
  \centering
  \includegraphics[width=0.9\linewidth]{resultados_pinn_2d.png}
  \caption{Validación PINN-SIREN para Helmholtz 2D con campo complejo.
           Fila superior: partes real e imaginaria de la solución analítica
           $E_{\mathrm{exacta}} = e^{i(k_x x + k_y y)}$.
           Fila inferior: predicción PINN-SIREN correspondiente.
           La diferencia visual es indistinguible; los errores puntuales son del orden
           de $10^{-3}$ (ver \cref{tab:resultados2d}).
           Malla de evaluación $100 \times 100$, SEED = 42.}
  \label{fig:resultados2d}
\end{figure}

\begin{figure}[htbp]
  \centering
  \includegraphics[width=\linewidth]{metricas_adicionales_2d.png}
  \caption{Métricas adicionales de entrenamiento NB02.
           Curva de pérdida (Adam + L-BFGS), mapa de error $|E_{\mathrm{real}} - \hat{E}_{\mathrm{real}}|$
           y $|E_{\mathrm{imag}} - \hat{E}_{\mathrm{imag}}|$ en la malla $100 \times 100$.
           Los errores más altos se concentran en las esquinas del dominio,
           región típicamente difícil para condiciones de onda plana.}
  \label{fig:metricas2d}
\end{figure}

\subsection{Generación de speckle óptico}
\label{subsec:resultados_speckle}

La \cref{tab:resultados_speckle} reporta la validación estadística del patrón de
speckle generado por la PINN-SIREN entrenada con condición de frontera rugosa
$\phi(x) \sim \mathcal{U}(0, 2\pi)$ en $y = 0$.

\begin{table}[htbp]
\centering
\begin{threeparttable}
\caption{Validación estadística del speckle --- PINN-SIREN NB03}
\label{tab:resultados_speckle}
\begin{tabular}{lcc}
\toprule
Métrica & Valor PINN & Valor teórico \\
\midrule
Contraste $C = \sigma_I / \langle I \rangle$ & $1.0253$ & $= 1$ \\
$|C - 1|$                                    & $0.0253$ & $< 0.1$ \\
KS $p$-valor                                 & $< 0.0001$ & $> 0.05$ \\
Frac.\ $I > 2\langle I \rangle$              & ${\approx}\,0.135$ & $e^{-2} \approx 0.1353$ \\
\midrule
Épocas Adam                                  & $7{,}976 / 15{,}000$ & --- \\
Iteraciones L-BFGS                           & $8 / 500$ & --- \\
Tiempo total (s)                             & $227$ & --- \\
\bottomrule
\end{tabular}
\begin{tablenotes}
\small
\item Evaluado en malla $100 \times 100$ puntos.
      Condición de frontera: $\phi(x) \sim \mathcal{U}(0, 2\pi)$, $N_\phi = 256$ pts.
      Bordes $x=0$, $x=1$, $y=1$: libres (sin condición Dirichlet artificial).
      Arquitectura y optimizador idénticos a NB02.
      GPU: NVIDIA RTX 5050, SEED = 42.
\end{tablenotes}
\end{threeparttable}
\end{table}

La condición de frontera $\phi(x) \sim \mathcal{U}(0, 2\pi)$, con cobertura completa del
círculo unitario, satisface la hipótesis de \citet{Goodman2007_Speckle} para speckle
completamente desarrollado.
El contraste $C = 1.0253$ cumple holgadamente el umbral $|C-1| < 0.1$ ($|C-1| = 0.025$),
y la fracción de píxeles con $I > 2\langle I \rangle$ es consistente con el valor teórico
$e^{-2} \approx 0.1353$, indicando que la distribución de colas es estadísticamente correcta.

El test KS falla formalmente ($p < 0.0001$, KS~$= 0.0628$).
Esta discrepancia se atribuye a la elevada potencia estadística del test con $N = 10{,}000$ puntos
y a la ausencia de condición de radiación de Sommerfeld en el borde $y = 1$, que en un dominio
abierto físico actuaría como frontera absorbente.
La incorporación de esta condición es parte del trabajo futuro (véase \cref{sec:conclusiones}).

La misma arquitectura SIREN $5\times128$ resuelve la ecuación de Helmholtz con cualquier
condición de frontera --- la transición de onda plana (NB02) a fase rugosa (NB03) no
requiere modificación alguna de la arquitectura, lo que demuestra la generalidad del
método.

\begin{figure}[htbp]
  \centering
  \includegraphics[width=0.48\linewidth]{resultados_speckle_nb03.png}
  \hfill
  \includegraphics[width=0.48\linewidth]{estadistica_speckle_nb03.png}
  \caption{Simulación de speckle óptico con PINN-SIREN (NB03).
           \textit{Izquierda:} Patrón de intensidad $I(x,y) = |E|^2$ en la malla
           $100 \times 100$, generado con condición de frontera
           $\phi(x) \sim \mathcal{U}(0, 2\pi)$ en $y=0$.
           El patrón granular es característico del speckle completamente desarrollado.
           \textit{Derecha:} Distribución de intensidades normalizada
           $p(I/\langle I\rangle)$ (histograma) comparada con la exponencial negativa
           teórica $e^{-I/\langle I\rangle}$ (línea roja).
           Contraste $C = 1.0253$ ($|C-1| = 0.025 < 0.1$).}
  \label{fig:speckle}
\end{figure}

\subsection{Comparativa con métodos de referencia}
\label{subsec:comparativa}

La \cref{tab:comparativa} resume el rendimiento de la PINN-SIREN propuesta
frente al estado del arte.

\begin{table}[htbp]
\centering
\begin{threeparttable}
\caption{Comparativa de error L2 --- PINN-SIREN vs estado del arte}
\label{tab:comparativa}
\begin{tabular}{lccc}
\toprule
Método & Dimensión & Error L2 (\%) & Factor de mejora \\
\midrule
\citet{Schoder2024_PINN_acoustics} & 1D & 2.490 & ---        \\
\citet{Schoder2024_PINN_acoustics} & 2D & 1.910 & ---        \\
\textbf{PINN-SIREN (este trabajo)} & \textbf{1D} & \textbf{0.006} & $\times 415$ \\
\textbf{PINN-SIREN (este trabajo)} & \textbf{2D} & \textbf{0.171} & $\times 11.2$   \\
\bottomrule
\end{tabular}
\begin{tablenotes}
\small
\item El factor de mejora es la razón entre el error de Schoder et al.\ (2024)
      y el error de la PINN-SIREN propuesta.
      Ambos trabajan con la ecuación de Helmholtz escalar en dominios normalizados.
      \citet{Schoder2024_PINN_acoustics} aborda el caso 3D con dominio físico y PINNs
      de activación \textit{tanh}; este trabajo opera en 1D/2D adimensionales con
      arquitectura SIREN.
      El factor de mejora refleja tanto la ventaja arquitectural de SIREN sobre
      \textit{tanh}/ReLU como la diferencia de complejidad dimensional.
\end{tablenotes}
\end{threeparttable}
\end{table}
